\documentclass[letterpaper,12pt]{article}
\usepackage{graphicx} % Required for inserting images

\usepackage[
  left=.9in,
  right=.9in,
  top=1in,
  bottom=1in,
  footskip=30pt
]{geometry}

\usepackage{csquotes}
\MakeOuterQuote{"}
\usepackage{comment}

\pagenumbering{gobble}

\title{Escaping Inequality in India: \\ The Role of English-Medium Instruction}

\author{Rishav Roy}
\date{}

\begin{document}

\maketitle
\vspace{-6mm}
\noindent What follows is my 49-page independent research paper, "Escaping Inequality in India: The Role of English-Medium Instruction." The paper's key points are described below.

\begin{enumerate}

\item \textbf{Question}: Does increasing baseline district-level EMI exposure (percentage of schoolgoing children enrolled in EMI) affect consumption growth over the next decade? 

\item \textbf{Data sources}: 2007-08/2017-18 microdata and metadata from the National Sample Survey, 35 datasets from the 2001 Census, and 2019-20 GIS shapefiles.
\begin{comment}
\vspace{-3mm}
\begin{itemize}
\item \textbf{Key variables}: Baseline EMI exposure (2007-08), consumption outcomes (2017-18), and linguistic distance instrument (2001).
\end{itemize}
\vspace{-2mm}
\end{comment}

\item \textbf{Identification}: District-level 2SLS (pop.-weighted linguistic distance from Hindi IV, state-clustered SEs). Individual-level probit of  education selection, design-based SEs. 
\vspace{-3mm}
\begin{itemize}
\item \textbf{Key clarification}: State FEs, essential for 2SLS exclusion restriction, trigger extreme multicollinearity. District matching \& possible FD-2SLS redesign needed.
\end{itemize}
\vspace{-2mm}

\item \textbf{Main result}: Strong first-stage ($F=39.2$), null second-stage estimate ($0.2$ pp, $p=0.78$). \underline{Estimates are provisional} pending district matching that enables state FEs.

\item \textbf{Next steps}: 
\vspace{-3mm}
\begin{itemize}
\item Repair 2SLS exclusion restriction with district harmonization changes, FD-2SLS.

\vspace{-1mm}

\item Redesign models for spatial spillovers, migration, `bad controls,' heterogeneity. 

\vspace{-1mm}

\item Make response variable robust to local variation in inflation and shocks.
\end{itemize}
\vspace{1mm}


\end{enumerate}


\noindent Due to the paper's length, I recommend the time-constrained reviewer prioritize these ten pages: 

\begin{enumerate}
\item[6:] \textbf{Research question}, contribution to literature, and summary/outline of paper.
\item[10-11:] \textbf{EMI exposure, probit model}, data wrangling, variable construction, missingness.
\item[23:] \textbf{Instrumental variable}, relevance, issues with exclusion restriction ($\kappa=28730.84$).
\item[25:] \textbf{Interpretation of 2SLS}, key limits to interpretation, planned remedies.
\item[26-27:] \textbf{2SLS results}, both first- and second-stage.
\item[28:] \textbf{Interpretation of 2SLS}, small \& noisy estimate, `bad controls,' spillover attenuation. 
\item[32-33:] \textbf{District harmonization}  methods, assumptions. Some next steps (planned regressors, interactions, potential Bartik instrument).

\end{enumerate}

\end{document}